本研究開發了一個針對澳門氹仔地區共享單車系統的綜合優化框架，涵蓋站點選址、初始投放量優化和再平衡問題。在選址與投放方面，模型通過將區域細分為居民區、旅遊區和樞紐區，並基於歷史需求數據估計淨需求流量，採用混合整數規劃方法確定最佳站點位置與初始投放量，以最小化高峰時段的服務缺口。結果顯示，21個站點和280輛單車的配置下，全局最大服務缺口為12，總缺口為65，滿足了空間容量與預算約束。在再平衡方面，模型構建了一個基於有向圖的調度方案，考慮運輸成本、懲罰成本及卡車運載量等實際約束，允許不完全滿足需求以提升靈活性，並支持多輛卡車調度。通過數值試驗，該模型在單車調度中實現了總成本的最優化（運輸成本15500，懲罰成本19000），為運營決策提供了參考。整體而言，該模型結合理論與實踐，為共享單車系統的高效運行提供了科學工具。